% sections/introduction.tex — Giriş bölümü, taslak v1 (2026-10-03)
% Kaynaklar: iskelet v6 (Giriş konu başlıkları), talimat §5; makale1–3.zip'teki tam metinler.
% Her sayı kaynağın kendi metninden: Raine1990 (s. 1003), Chan2023 (s. 6), Berluti2026 (s. 390, 395),
% Tasci2025 (özet ve s. 15), Bonfante2023 (Tablo), Steele2017 (s. 7–8), Gil2026 (özet), Varoquaux2018 (özet).

\section{Introduction}
\label{sec:intro}

The search for biological markers of antisocial behaviour has a long history. One of its oldest hypotheses is low arousal: in a prospective study, adolescents who were later convicted had shown a lower resting heart rate, lower skin conductance and more slow-wave EEG activity than those who were not \cite{Raine1990}. Machine learning now promises to turn such group differences into individual classifications, and recent studies report high performance. A neural network predicted conduct disorder two years ahead with 91\% accuracy in a single random hold-out split \cite{Chan2023}, and random forests predicted later conduct problems with AUCs of 0.97--0.98 from predictors that included the children's baseline diagnoses \cite{Berluti2026}. From EEG, `psychotic criminals' were distinguished from controls with up to 99.95\% accuracy \cite{Tasci2025}. A study from the research group that collected the present data separated adolescent offenders from non-offenders with cognitive test scores, with a balanced accuracy of 0.88 \cite{Bonfante2023}.

High performance does not by itself show a marker that generalises. Three problems are common. The first is data leakage. When segments of the same participant appear in both training and test sets, a classifier can recognise participants instead of groups \cite{Saeb2017,Brookshire2024}; in EEG, performance estimated across time windows can fall sharply once whole subjects are held out \cite{Gil2026}. Selecting features or tuning a model on the data used for evaluation also inflates the estimates \cite{Varma2006,Cawley2010,Kapoor2023}. The second is sample size: with about 100 participants, cross-validated estimates carry error bars of roughly $\pm 10$ percentage points \cite{Varoquaux2018}. The third is confounding. When groups are recruited, recorded or scanned in different places or periods, a model may capture the circumstances of the measurement rather than the condition of interest \cite{Chyzhyk2022}. In a structural MRI study of incarcerated adolescents, offenders were scanned at the detention facility and controls at a research centre, and age and IQ alone separated offenders from controls at least as well as grey-matter features \cite{Steele2017}. A recent scoping review of machine learning in schizophrenia-spectrum disorders found that most studies did not address generalisability in substance \cite{Rivera2026}.

The open dataset analysed here (OpenNeuro ds006923) is a test case for all three problems \cite{Polo2025ds006923}. It contains resting-state EEG (128 channels) of 140 male Colombian adolescents: 74 juvenile offenders, labelled as two subgroups (sg and sg2), and 66 controls. The recordings come from a project that also studied interoception in these participants \cite{RodriguezAlvarez2023}, and the same research group has studied cognition, emotion processing and brain structure in young offenders \cite{Bonfante2023,SanchezTorres2024,RuizPena2024,Pino2023,Pino2019,FrancoOByrne2021,SantamariaGarcia2019}. The dataset invites classification: each participant contributes several segments, and precomputed spectral features are provided for most participants. Its design, however, makes group differences hard to attribute. Group is confounded with recording period: sg was recorded in 2021 and never in the same month as the controls, whereas sg2 overlapped with the controls in only two months. Offenders and controls also differ in substance use, schooling, socioeconomic background and age (table~\ref{tab:Table1}).

Leakage from segment-wise splitting is well documented in EEG \cite{Brookshire2024,Gil2026}. Our contribution is to audit an open forensic EEG dataset as a whole: the identity of its participants, the way its distributed features were computed, the confounding of group with recording period, and the transfer of a separation across offender subgroups. We asked how much of the apparent separability of offenders and controls survives correct validation, and whether any remaining separation generalises. We (i) audited the recordings and the precomputed features; (ii) quantified leakage on the precomputed features by comparing naive with nested, subject-wise pipelines; (iii) tested a single hypothesis that was designated as primary before it was tested, although after an initial analysis of the same data (table~\ref{tab:TableLog}): a group difference in alpha reactivity to eye opening, an index related to arousal \cite{Barry2007} and thus to the underarousal hypothesis \cite{Raine1990}; (iv) quantified the null result of prespecified region-of-interest features; and (v) in exploratory analyses, tested whether a channel-level separation of one offender subgroup from the controls transfers to the second subgroup, which is a direct test of generalisation. Every plan, decision rule and change was recorded in a time-stamped analysis log.
